\documentclass[11pt,a4paper]{article}

% =========================================================================
% PAKET-PAKET UTAMA
% =========================================================================
\usepackage[utf8]{inputenc}
\usepackage[T1]{fontenc}
\usepackage{lmodern}
\usepackage[indonesian]{babel}
\usepackage{amsmath,amssymb,amsfonts,amsthm}
\usepackage{bm}              % Notasi tebal untuk vektor & matriks
\usepackage{geometry}
\geometry{top=1.8cm, bottom=1.8cm, left=2.0cm, right=2.0cm}
\usepackage{graphicx}
\usepackage{booktabs}        % Tabel profesional
\usepackage{tabularx}
\usepackage{microtype}
\usepackage{xcolor}
\usepackage{mdframed}        % Kotak highlight untuk definisi
\usepackage{caption}
\usepackage{subcaption}
\usepackage{float}
\usepackage{tikz}
\usetikzlibrary{shapes.geometric, arrows.meta, positioning, calc}
\usepackage{hyperref}
\hypersetup{
    colorlinks=true,
    linkcolor=blue!70!black,
    citecolor=blue!70!black,
    urlcolor=blue!70!black,
    pdfencoding=auto,
    psdextra
}
\pdfstringdefDisableCommands{%
  \def\boldsymbol#1{#1}%
  \def\bm#1{#1}%
  \def\mathbf#1{#1}%
}

% Makro sitasi fleksibel
\newcommand{\citep}[2][]{%
  \def\tempa{#2}%
  \def\tempb{rasmussen2006gaussian}%
  \def\tempc{damianou2013deep}%
  \def\tempd{salimbeni2017doubly}%
  \def\tempe{damianou2013deep,salimbeni2017doubly}%
  \ifx\tempa\tempb(Rasmussen \& Williams, 2006\if\relax\detokenize{#1}\relax\else, #1\fi)%
  \else\ifx\tempa\tempe(Damianou \& Lawrence, 2013; Salimbeni \& Deisenroth, 2017)%
  \else\ifx\tempa\tempc(Damianou \& Lawrence, 2013)%
  \else\ifx\tempa\tempd(Salimbeni \& Deisenroth, 2017)%
  \else(#2\if\relax\detokenize{#1}\relax\else, #1\fi)%
  \fi\fi\fi\fi%
}
\newcommand{\citet}[2][]{Rasmussen \& Williams (2006\if\relax\detokenize{#1}\relax\else, #1\fi)}

% Lingkungan Teorema & Definisi
\theoremstyle{plain}
\newtheorem{teorema}{Teorema}[section]
\newtheorem{claim}{Claim}[section]

\theoremstyle{definition}
\newtheorem{definisi}{Definisi}[section]

\theoremstyle{remark}
\newtheorem{catatan}{Catatan}[section]

% Definisi warna kustom untuk kotak dan diagram
\definecolor{dgpblue}{RGB}{26, 82, 118}
\definecolor{dgpgreen}{RGB}{30, 132, 73}
\definecolor{dgppurple}{RGB}{108, 52, 131}
\definecolor{dgpred}{RGB}{148, 49, 38}

\newmdenv[
    linecolor=dgpblue,
    outerlinewidth=1.2pt,
    roundcorner=4pt,
    innertopmargin=6pt,
    innerbottommargin=6pt,
    innerrightmargin=10pt,
    innerleftmargin=10pt,
    backgroundcolor=blue!2
]{mathbox}

\newmdenv[
    linecolor=dgpred,
    outerlinewidth=1.2pt,
    roundcorner=4pt,
    innertopmargin=6pt,
    innerbottommargin=6pt,
    innerrightmargin=10pt,
    innerleftmargin=10pt,
    backgroundcolor=red!2
]{alertbox}

% Format nama pembuktian
\addto\captionsindonesian{\renewcommand{\proofname}{Proof}}
\renewcommand{\proofname}{Proof}

% Path direktori gambar
\graphicspath{{figures/}}

% =========================================================================
% INFORMASI DOKUMEN
% =========================================================================
\begin{document}

\begin{center}
    {\Large\textbf{BIMBINGAN TUGAS AKHIR \#7}} \\[0.3em]
    {\small Selasa, 27 Oktober 2026}
\end{center}
\vspace{-0.2cm}
\setlength{\abovedisplayskip}{6pt}
\setlength{\belowdisplayskip}{6pt}
\setlength{\abovedisplayshortskip}{4pt}
\setlength{\belowdisplayshortskip}{4pt}

% =========================================================================
% 1. p(y|X) PADA GP DAN DGP
% =========================================================================
\section{\texorpdfstring{$p(\mathbf{y} \mid X)$}{p(y|X)} pada GP dan DGP}
\label{sec:py_gp_dgp}

% -------------------------------------------------------------------------
\subsection{Distribusi Marginal \texorpdfstring{$p(\mathbf{y} \mid X)$}{p(y|X)} pada Gaussian Process}
\label{subsec:marginal_gp}

\noindent Ditinjau model $\mathbf{y} = X \mathbf{w} + \boldsymbol{\varepsilon}$ dengan $\mathbf{w} \sim \mathcal{N}\left(\mathbf{0}, \; \Sigma_p\right)$ dan $\boldsymbol{\varepsilon} \sim \mathcal{N}\left(\mathbf{0}, \; \sigma_n^2 I_N\right)$:
\begin{equation}
    \mathbf{y} \mid X \sim \mathcal{N}\left(\mathbf{0}, \; X \Sigma_p X^T + \sigma_n^2 I_N\right) \label{eq:marginal_linier}
\end{equation}
dengan kovariansi linier $K_{\text{lin}} = X \Sigma_p X^T$.

\noindent Untuk memperluas model ke fungsi non-linier, vektor masukan $\mathbf{x} \in \mathbb{R}^D$ dipetakan ke dalam ruang fitur (\textit{feature space}) $\mathcal{H}$ berdimensi tinggi melalui pemetaan $\boldsymbol{\phi}(\mathbf{x}): \mathbb{R}^D \to \mathcal{H}$ \citep[hlm.~11--12]{rasmussen2006gaussian}. Matriks masukan $X \in \mathbb{R}^{N \times D}$ bertransformasi menjadi matriks fitur $\Phi = [\boldsymbol{\phi}(\mathbf{x}_1), \dots, \boldsymbol{\phi}(\mathbf{x}_N)]^T \in \mathbb{R}^{N \times |\mathcal{H}|}$, sehingga struktur kovariansi $X \Sigma_p X^T$ bertransformasi menjadi:
\begin{equation}
    K = \Phi \Sigma_p \Phi^T \in \mathbb{R}^{N \times N} \label{eq:kovariansi_fitur}
\end{equation}
Melalui penerapan \textit{kernel trick}, setiap elemen matriks kovariansi $K$ dievaluasi secara langsung melalui fungsi kernel tanpa komputasi eksplisit di ruang $\mathcal{H}$:
\begin{equation}
    K_{ij} = k(\mathbf{x}_i, \mathbf{x}_j) = \boldsymbol{\phi}(\mathbf{x}_i)^T \Sigma_p \boldsymbol{\phi}(\mathbf{x}_j) \label{eq:kernel_trick}
\end{equation}
Dengan demikian, distribusi marginal observasi pada Gaussian Process adalah:
\begin{equation}
    \mathbf{y} \mid X \sim \mathcal{N}\left(\mathbf{0}, \; K + \sigma_n^2 I_N\right) \label{eq:distribusi_y_gp}
\end{equation}

\noindent Fungsi kerapatan probabilitas marginal (\textit{marginal likelihood}) $p(\mathbf{y} \mid X)$ dievaluasi dalam bentuk analitik tertutup (\textit{closed-form}):
\begin{equation}
    p(\mathbf{y} \mid X) = \frac{1}{(2\pi)^{N/2} \left| K + \sigma_n^2 I_N \right|^{1/2}} \exp\left( -\frac{1}{2}\mathbf{y}^T \left(K + \sigma_n^2 I_N\right)^{-1} \mathbf{y} \right) \label{eq:marginal_likelihood_densitas}
\end{equation}
dengan fungsi \textit{log marginal likelihood}:
\begin{equation}
    \log p(\mathbf{y} \mid X) = -\frac{1}{2}\mathbf{y}^T \left(K + \sigma_n^2 I_N\right)^{-1} \mathbf{y} - \frac{1}{2}\log \det\left(K + \sigma_n^2 I_N\right) - \frac{N}{2}\log(2\pi) \label{eq:log_marginal_likelihood}
\end{equation}

% -------------------------------------------------------------------------
\subsection{Bentuk \texorpdfstring{$p(\mathbf{y} \mid X)$}{p(y|X)} pada Deep GP}
\label{subsec:intraktabilitas_dgp}

\noindent Pada arsitektur \textit{Deep Gaussian Process} (DGP) dengan $L$ lapisan hierarkis, pemetaan dimodelkan secara berantai:
\begin{align}
    \mathbf{h}^{(0)} &= X \in \mathbb{R}^{N \times D_{\text{in}}} \label{eq:dgp_input} \\
    \mathbf{h}^{(\ell)} &= \mathbf{f}^{(\ell)}\left(\mathbf{h}^{(\ell-1)}\right) + \boldsymbol{\varepsilon}^{(\ell)}, \quad \mathbf{f}^{(\ell)} \sim \mathcal{GP}\left(\mathbf{0}, \; k^{(\ell)}\left(\mathbf{h}^{(\ell-1)}, \mathbf{h}^{(\ell-1)'}\right)\right), \quad \ell = 1, \dots, L-1 \label{eq:dgp_hidden} \\
    \mathbf{y} &= \mathbf{f}^{(L)}\left(\mathbf{h}^{(L-1)}\right) + \boldsymbol{\varepsilon}^{(L)}, \quad \boldsymbol{\varepsilon}^{(\ell)} \sim \mathcal{N}\left(\mathbf{0}, \; \sigma_\ell^2 I_N\right) \label{eq:dgp_output}
\end{align}

\noindent Berdasarkan pemodelan lapisan pada Persamaan~\eqref{eq:dgp_hidden} dan \eqref{eq:dgp_output}, distribusi marginal bersama $p(\mathbf{y} \mid X)$ dituliskan dengan mengintegrasikan seluruh variabel laten $\mathbf{h}^{(1)}, \dots, \mathbf{h}^{(L-1)}$:
\begin{equation}
    p(\mathbf{y} \mid X) = \int p\left(\mathbf{y} \mid \mathbf{h}^{(L-1)}\right) \left[ \prod_{\ell=1}^{L-1} p\left(\mathbf{h}^{(\ell)} \mid \mathbf{h}^{(\ell-1)}\right) \right] d\mathbf{h}^{(1)} \dots d\mathbf{h}^{(L-1)}
    \label{eq:integral_dgp_marginal}
\end{equation}

\end{document}
